\section{Executed experiments}
\label{sec:experiments}

\subsection{Thirty-six runs, reported separately}

This section reports thirty-six independent prototype runs in four rounds: the
design round \src{p1}--\src{p9}, covered first; the extension round
\src{e1}--\src{e9} (Section~\ref{sec:extension-runs}); the third round
\src{r1}--\src{r9} (Section~\ref{sec:round-three-runs}); and the fourth
\src{s1}--\src{s9} (Section~\ref{sec:round-four-runs}). All thirty-six are
named after the source proposals listed in Appendix~\ref{app:provenance}. They
are reported separately and they are \emph{not} pooled --- neither within a
round nor across the four --- for a reason stronger than caution.

All nine of the design-round runs ran on the same day in the same kind of
environment --- Python 3.13.5
on Linux x86-64, NumPy 2.3.5, SciPy 1.17.0, SymPy 1.14.0, and where pytest was
used, pytest 9.0.2 --- and five of them recorded the same seed value,
\code{20260914}. That coincidence makes the runs look comparable. They are not.
The same seed drives entirely different generators over entirely different case
sets: one run generates 30 cone and 60 quadratic cases, another 1{,}080 witness
tables and 300 CDCL(T) instances, another 100 cone and 49 invariant cases,
another 1{,}500 differential SAT cases, another 26 curated certificates. Three
runs (\src{p1}, \src{p8}, \src{p9}) record no seed at all.

\begin{remark}[The non-pooling rule]
The counts below are not commensurable. ``Cases'', ``checks'', ``assertions'',
``tests'', ``certificates'' and ``bundles'' are five different units, and several
runs use more than one. Summing them would produce a meaningless total, and any
such total would describe the prototypes' own constructed distributions rather
than any theorem corpus. Every number in this section carries its run label for
that reason.
\end{remark}

Two properties do hold uniformly, and they are the ones worth carrying forward.
First, no coefficient anywhere in the certificate corpus is stored as a
floating-point number: rationals are exact throughout, serialised as strings
such as \code{"125/18"} or as integer numerator/denominator pairs. A scan of
every certificate payload in all nine runs found exactly one float, and it is a
timing field. Second, every run's stored certificates replay under a
standard-library-only checker, verified with site packages disabled.

\subsection{The runs}

\subsubsection*{\src{p1} --- six certificate families, one uniform protocol}

275 pytest instances passed. A deterministic benchmark of 198 constructed
instances produced certificates in every case, and a separate
standard-library-only decoder accepted 198/198. Each case was searched three
times after numerical warm-up; the table reports medians of per-case medians.

\begin{table}[htbp]
\centering\small
\begin{tabular}{@{}lrrrr@{}}
\toprule
Component & Certified & Ablation & Search (ms) & Recheck (ms) \\
\midrule
Quadratic SOS            & 60/60 & 1 & 0.864 & 0.701 \\
Finite cone LP           & 17/17 & 0 & 1.699 & 0.295 \\
Bernstein box            & 18/18 & 1 & 0.685 & 0.389 \\
Polynomial recurrence    & 33/33 & --- & 2.513 & 0.403 \\
Integral affine witness  & 40/40 & --- & 0.074 & 0.027 \\
Univariate with zeros    & 30/30 & 0 & 1.930 & 0.505 \\
\bottomrule
\end{tabular}
\caption{\src{p1}. Ablations are restricted versions of these same prototype
mechanisms, \textbf{not} \grind{} and not any Lean tactic. A dash means no
ablation was run. Search time includes checking performed internally by the
search routine, so the recheck column is not the amount of checking hidden
inside the search column.}
\label{tab:p1}
\end{table}

The ablations: the quadratic one recognises only expanded nonnegative
combinations of even monomials; the cone one removes richer square candidates
or, for the domain example, the hypothesis products; the Bernstein one disallows
subdivision; the univariate one runs dyadic Bernstein search to depth 12 without
factorisation. The last row is the sharpest result in this run --- 30/30 against
0/30 --- and it is the experimental face of Theorem~\ref{thm:zero}.

Mutation tests alter targets, weights, hypotheses, recurrence base terms, affine
offsets, Bernstein split points, child nodes, claimed bounds, and Sturm chains,
and require rejection. Independent symbolic tests reconstruct Bernstein
expansions, expand quadratic certificates, and compare univariate root counts.
Limitations are preserved as regression cases: a degree-limited recurrence
search failing on a valid higher-degree answer, integer affine synthesis
rejecting a merely rational witness, Bernstein-only search failing at a
non-dyadic interior zero.

\subsubsection*{\src{p2} --- five families, per-family seeds}

31 unittest methods passed; several contain multiple seeded trials. 215
certificates were produced and all 215 replayed from storage.

\begin{table}[htbp]
\centering\small
\begin{tabular}{@{}lrrr@{}}
\toprule
Family & Cases & Exact replays & Time (s) \\
\midrule
List induction and generalisation & 6 & 6 & 0.005 \\
Polynomial additive recurrences & 65 & 65 & 0.241 \\
Affine witnesses with Farkas certificates & 63 & 63 & 0.352 \\
Finite-dictionary square certificates & 40 & 40 & 0.648 \\
Bernstein box certificates & 41 & 41 & 0.067 \\
\midrule
Total & 215 & 215 & \\
\bottomrule
\end{tabular}
\caption{\src{p2}. Seeds are per family: 420 recurrences, 421 witnesses,
422 squares, 423 Bernstein. Times are one local run and are not
hardware-normalised.}
\label{tab:p2}
\end{table}

The recurrence family comprises five named power sums and 60 seeded random
polynomial increments; it required 312 candidates in total, and the exact step
checks generated 247 counterexamples before the 65 certificates were accepted.
The 41 Bernstein problems used 82 leaves in total at maximum split depth four;
root-only checking certifies 25 of them and subdivision certifies all 41. The
six list proof traces contain 6, 3, 8, 10, 9 and 5 rewrites.

The run's own caveat is the important one: \emph{the 40 square-certificate
problems are constructed from the same finite dictionary the search uses}, so
their 100\% success rate tests implementation and rational reconstruction on a
designed distribution and must not be read as nonlinear proving coverage.
Recorded negative outcomes include the Motzkin polynomial finding no certificate
in the plain-square dictionary, the constant witness template failing on the
unbounded strip, and $(x-y)^2$ remaining \unknown{} at depth four.

\subsubsection*{\src{p3} --- assertion-counted mechanism ablations}

1{,}680 assertions passed; 24 certificate bundles were saved and all 24 replayed
under a standard-library-only checker. Seed \code{20260914}; total 2.27 s.

\begin{table}[htbp]
\centering\small
\begin{tabular}{@{}lrr@{}}
\toprule
Mechanism & Restricted & Enabled \\
\midrule
Cone certificates, curated goals & 0/10 diagonal-only & 9/10 full dictionary \\
Cone certificates, constructed & --- & 30/30 \\
Exact quadratic, constructed & --- & 60/60 \\
Exact quadratic, curated & --- & 4/4 \\
Bernstein box, 10 true inequalities & 3/10 unsplit & 8/10 adaptive \\
Ground Horn, 12 goals & 3/12 unsliced & 12/12 demand-sliced \\
List induction, 8 goals & 6/8 fixed accumulator & 8/8 generalised \\
\bottomrule
\end{tabular}
\caption{\src{p3}. Every ``restricted'' column is a restricted version of this
run's own implementation. None is \grind{}.}
\label{tab:p3}
\end{table}

Four negative controls were rejected by the quadratic path, seven corruptions by
the Bernstein checker, and eighteen by the induction checker. 250 random Horn
graphs and 500 comparisons confirmed that the sliced and unsliced modes agree
with a naive least-fixed-point computation. Both Horn modes still pay the full
index-construction cost, and the examples deliberately place decoys first.

The 60 generated quadratic examples are constructed from positive square
combinations and therefore carry an explicit selection bias; the ten cone goals
are deliberately selected mechanism tests, not an unbiased corpus.

\subsubsection*{\src{p4} --- the largest artefact corpus}

64 pytest tests passed. The deterministic experiment produced 1{,}494 rows and
1{,}238 saved successful artefacts, all replayed by checking modules that import
no numerical or symbolic search library. Seed \code{20260914}. This run is also
the only one that re-verified its corpus from an extracted archive copy and
confirmed that every family and status count matched the recorded run.

\begin{table}[htbp]
\centering\small
\begin{tabular}{@{}lr@{}}
\toprule
Family & Observed result \\
\midrule
Residue witness parameters & 1{,}080 checked: 831 proved, 249 impossible by gcd divisibility \\
Random CNF + integer difference logic & 300 oracle agreements: 189 UNSAT, 111 SAT \\
Sparse cone / equality-ideal certificates & 47 of 49; two deliberate \unknown{} \\
Exact rational nonnegative quadratics & 40/40 manufactured examples \\
Bernstein covering certificates & 12 of 13; the non-dyadic zero case \unknown{} \\
Pigeonhole CNF & 5 replayed UNSAT proofs \\
Structural equations & 3 synthesised or proved lemmas, 2 rejected false candidates \\
Induction ablation & 4 rows: 1 proved, 3 \unknown{} \\
\bottomrule
\end{tabular}
\caption{\src{p4}. These families are heterogeneous and must not be combined
into a general accuracy score.}
\label{tab:p4}
\end{table}

The structural ablation is this run's sharpest finding, and it is a genuine
$2\times2$ design rather than a single comparison:

\begin{table}[htbp]
\centering\small
\begin{tabular}{@{}ccc@{}}
\toprule
Proved append lemmas & Generalise accumulator & Outcome \\
\midrule
No  & No  & \unknown{} \\
No  & Yes & \unknown{} \\
Yes & No  & \unknown{} \\
Yes & Yes & \textsc{proved} \\
\bottomrule
\end{tabular}
\caption{\src{p4}. Only generalisation \emph{and} proved auxiliary lemmas
together close the accumulator reversal. The same goal and the same rewriting
engine are used in all four configurations, which is stronger evidence of
mechanism than a single all-features-enabled demonstration, and it shows why a
blanket ``try induction'' addition would miss part of the design.}
\label{tab:p4-ablation}
\end{table}

Two further details behind the aggregate counts. The 1{,}080 residue cases sweep
all $1\le m\le20$, $-3\le a\le5$ and $0\le b\le5$; the 249 impossible ones fail
the exact gcd divisibility condition. And for
$p(x)=(x-1/3)^2+10^{-4}$ on $[0,1]$ --- the interior-zero example given a
positive margin --- the prototype returns a tree with seven leaves, thirteen
nodes and maximum depth six, whereas $(x-1/3)^2$ itself reaches the depth limit,
exactly as Theorem~\ref{thm:zero} requires.

The run's own caveats are worth preserving: 64 is a test-runner count, not a
count of mathematical instances, since parametrised tests contain inner-loop
cases; forty of the 49 cone cases are manufactured from the search's own
certificate basis; the recorded timing is one in-process pass rather than a
repeated performance study; and the field named \code{certificate\_bytes}
includes the accompanying statement, not a proof payload alone. The run states
the non-pooling rule for itself in the sharpest available form: \emph{do not
divide 1{,}238 by 1{,}494 and advertise an overall accuracy}, because the
denominator mixes proof searches, model searches, a lemma batch, deliberately
false formulas, and overlapping ablation runs.

The pigeonhole family is recorded per instance, including the datum that is
least flattering to the method:

\begin{table}[htbp]
\centering\small
\begin{tabular}{@{}rrrrr@{}}
\toprule
Pigeons / holes & Decisions & Conflicts & Resolution nodes & Nonchronological jumps \\
\midrule
2 / 1 &  0 &  1 &   2 & 0 \\
3 / 2 &  1 &  2 &  10 & 0 \\
4 / 3 &  8 &  9 &  45 & 0 \\
5 / 4 & 30 & 31 & 160 & 0 \\
6 / 5 & 89 & 89 & 513 & 1 \\
\bottomrule
\end{tabular}
\caption{\src{p4}. The small number of non-chronological jumps is reported
rather than hidden. It is not evidence that this branch heuristic is competitive
on industrial SAT instances. Note also that this encoding --- at least one hole
per pigeon and at most one pigeon per hole --- differs from \src{p6}'s, so
Tables~\ref{tab:p6-dpll} and this one are not comparable.}
\label{tab:p4-pigeonhole}
\end{table}

\subsubsection*{\src{p5} --- expected-outcome benchmarking with negative controls}

73 pytest tests passed in 1.11 s. Seed \code{20260914}.

\begin{table}[htbp]
\centering\small
\begin{tabular}{@{}lll@{}}
\toprule
Experiment & Outcome & Interpretation \\
\midrule
Named nonlinear cases & 9 accepted / 12 &
Two true cases unresolved; one deliberately false case unresolved \\
Generated product-cone cases & 100 / 100 &
Constructed inside the covered cone; not representative sampling \\
False-polynomial controls & 0 accepted / 24 &
All \unknown{}; exact negative witnesses checked separately \\
Bernstein cases & 43 / 45 &
One true non-dyadic-zero case and one false case unresolved \\
Additive recurrence invariants & 49 / 49 &
Base and universal step identities checked \\
Horn growth family & 5 depths &
Both strategies prove all goals; search work differs sharply \\
Actual Lean comparison & Not run & No measured superiority over \grind{} \\
\bottomrule
\end{tabular}
\caption{\src{p5}.}
\label{tab:p5}
\end{table}

The run states its own most important qualification: the \textbf{109} accepted
nonlinear certificates are the nine named successes plus the 100 generated
cases, and they are \emph{not} 109 previously unsolved Lean theorems. The 24
negative controls are of the form $(x-a)^2-\varepsilon$ with rational
$\varepsilon>0$; the benchmark separately checks the exact negative value at $a$,
while the cone routine itself returns \unknown{} rather than a certified
refutation. Median recorded cone discovery time was about 21.5\,ms with a
maximum of about 103.6\,ms; median invariant synthesis about 4.1\,ms, including
interpolation and validation but excluding any Lean work. 58 exported
certificates --- nine named cone plus all 49 invariants --- were reloaded and
re-accepted by a separate validator.

\subsubsection*{\src{p6} --- proof-logging CDCL(T)}

This run's core is a differential test against an exhaustive reference.

\begin{table}[htbp]
\centering\small
\begin{tabular}{@{}lr@{}}
\toprule
Check & Result \\
\midrule
Random pure CNF cases & 1{,}000 \\
Random CNF + integer-difference cases & 500 \\
Satisfiable results & 1{,}278 \\
Unsatisfiable results & 222 \\
Differential mismatches & 0 \\
UNSAT certificates independently replayed & 222 \\
Truncated certificates rejected & 222 \\
Additional structural edge cases & 5 \\
Deliberately mutated theory certificates rejected & 2 \\
\bottomrule
\end{tabular}
\caption{\src{p6}. Formulas have three to eight Boolean variables, clauses of
three distinct variables with independent signs, and one to six times the
variable count in clauses; difference constraints use four integer variables and
constants in $[-3,3]$. The reference enumerates all Boolean assignments and uses
an all-pairs shortest-path consistency check; the solver under test uses CDCL
with negative-cycle theory conflicts. These are small exhaustive tests, not an
industrial benchmark.}
\label{tab:p6-sat}
\end{table}

The five edge cases are an empty conjunction, an empty clause, a tautology,
contradictory units, and repeated literals; the integer-complement boundary has
its own regression assertion.

\begin{table}[htbp]
\centering\small
\begin{tabular}{@{}lrrrr@{}}
\toprule
Instance & Variables & Clauses & DPLL decisions & CDCL decisions \\
\midrule
Padded core, 0 unused  &  2 &  4 &     1 &   1 \\
Padded core, 4 unused  &  6 &  4 &    31 &   5 \\
Padded core, 8 unused  & 10 &  4 &   511 &   9 \\
Padded core, 12 unused & 14 &  4 & 8{,}191 &  13 \\
Pigeonhole $4\to3$     & 12 & 22 &     8 &   8 \\
Pigeonhole $5\to4$     & 20 & 45 &    51 &  39 \\
Pigeonhole $6\to5$     & 30 & 81 &   374 & 148 \\
\bottomrule
\end{tabular}
\caption{\src{p6}. The baseline is \emph{this run's own} chronological
unit-propagating DPLL with the same static variable order and phase. It is not
\grind{} and not a production SAT solver.}
\label{tab:p6-dpll}
\end{table}

The run attaches the caveats itself, and they should be kept. The padded-core
family is an artificial diagnostic: it places unused Boolean variables before a
two-variable inconsistent core, which a production preprocessor would delete.
Its decision-count gap illustrates non-chronological backjumping; it is not
credible evidence of practical speedup. More pointedly, \emph{runtime gives a
soberer picture than decision counts}: on the largest pigeonhole case both
solvers took about 26\,ms despite their different decision counts, with replay
adding about 21\,ms. Operation counts are treated as diagnostic and no general
speed claim is made.

Ten polynomial search cases yielded seven certificates and three deliberate
\unknown{}s --- a false target, Motzkin at degree cap 6, and $(2x+y)^2$ at degree
cap 2. A further 100 randomly planted rational-square certificates exercised the
checker and were cross-checked symbolically; those are \emph{checker tests, not
100 additional search successes}, and the manually supplied discriminant
identity \eqref{eq:disc} is likewise separate from the seven discoveries. Six
invariant powers $p=0,\ldots,5$ yielded degrees 1 to 6 with final sample counts
3, 4, 5, 6, 7, 8, and each rejected the alterations $Q+1$ and $Q+n$, giving 12
mutation rejections on the synthesis side and 12 more under independent replay.

The saved corpus is 8 SAT/theory proof logs, 8 polynomial certificates (seven
discovered, one manually supplied), and 6 invariant certificates, all replayed
with site packages disabled. Three polynomial search files carry no certificate
and are explicitly skipped as \unknown{}. The replay program fails on a missing
corpus rather than reporting vacuous success.

\subsubsection*{\src{p7} --- component cases with corruption controls}

833 principal component test cases produced 891 accepted certificate checks;
330 deliberately corrupted certificates were rejected; 150 exact polynomial
cross-checks were run against a symbolic oracle; and 897 synthesised accumulator
formulas were compared against recursive execution. Seed \code{20260914}.

The run is explicit that \emph{cases and checks are different units}: some Horn
cases compare two successful derivations while others have no derivation, and
timing repetitions and archived demonstrations are excluded from the aggregate
check count. Five cone-search misses and two invariant-template misses are
expected outcomes, not wrong proofs.

26 representative certificates --- 9 cone identities, 5 lattice traces, 11
accumulator invariants, 1 Horn derivation --- all passed separate
standard-library-only replay. The lattice worker was additionally validated by
an exhaustive 539-case single-row grid checked against integer gcd, and by 120
random feasible systems whose free-basis samples were re-checked against every
row.

\subsubsection*{\src{p8} --- the cleanest dictionary ablation}

260 pytest tests passed, the highest count in the design round. No seed
recorded.

\begin{table}[htbp]
\centering\small
\begin{tabular}{@{}lrr@{}}
\toprule
Experiment & Restricted baseline & Extended mechanism \\
\midrule
Polynomial cones, 60 constructed instances & 17 certificates & 60 certificates \\
Positive rational-box instances, 21 & 0 unsplit & 21 subdivided \\
Affine accumulator templates, 121 & not measured & 121 synthesised and replayed \\
Out-of-template recurrences, 5 & not applicable & 5 \unknown{} outcomes \\
\bottomrule
\end{tabular}
\caption{\src{p8}. The cone row varies only the dictionary and holds the LP
backend fixed, which makes it the most cleanly isolated ablation in the
collection. It is still an ablation of this run's own implementation.}
\label{tab:p8}
\end{table}

The Horn workload is the largest: a useful chain of length 12 buried in 1{,}000
unrelated chains of length 24, each with an initial fact. Full closure fires all
24{,}012 rules; an indexed baseline that stops as soon as the goal is known still
fires 11{,}012 under the specified interleaved queue order; backward slicing
selects just the 12 useful rules.

The accompanying observation matters more than the ratio: \emph{proof
minimisation reduces even the full-closure derivation to the same 12 proof
steps}, so proof size alone would conceal the entire search-work difference.
Conversely the cold reverse index still scans all input rules and the demand
implementation still loads the initial fact set, so no sublinear end-to-end cold
performance is claimed. The run explicitly declines to infer a speedup factor,
noting a \emph{slower} early-goal run despite fewer firings.

Supporting evidence: 121 invariant cases each cross-checked against runtime
execution with four deterministic random-list checks per pair (484 concrete
recurrence checks); 30 seeded sparse-polynomial tests against a symbolic oracle;
40 randomised Horn programs comparing full and sliced success and replaying
their derivations; and monomial-enumeration counts checked against
$\binom{n+d}{d}$. Only two certificates were exported to files, the fewest of
the nine.

\subsubsection*{\src{p9} --- small, heavily mutation-tested}

69 pytest tests passed. 13 persisted bundles replayed with site initialisation
disabled: one ordered bank of five inductively proved lemmas with a
specialisation proof, four Horn proofs, six cone certificates, one box
certificate, and one witness with its two side-condition proofs. These are
bundles, not 13 independent Lean theorems. No seed recorded.

\begin{table}[htbp]
\centering\small
\begin{tabular}{@{}lrrrr@{}}
\toprule
Depth & Forward facts & Backward proof nodes & Forward cap & Forward status \\
\midrule
 4 &    16 &  5 & 5{,}000 & Certified \\
 8 &   256 &  9 & 5{,}000 & Certified \\
12 & 4{,}096 & 13 & 5{,}000 & Certified \\
16 & 5{,}000 & 17 & 5{,}000 & Cap reached \\
\bottomrule
\end{tabular}
\caption{\src{p9}. Neither column is an implementation of \grind{}. The forward
baseline stops as soon as the target is produced, which is why its counts are
$2^d$ rather than $2^{d+1}-1$.}
\label{tab:p9-horn}
\end{table}

Mechanism outcomes: accumulator-lemma discovery enumerates 41 candidates,
filters against 49 sample environments, and the 41st enumerated term
$\mathrm{app}(\mathrm{rev}(xs),a)$ is the one that receives an induction
certificate; the quartic nonnegativity case builds 280 candidate columns and 35
coefficient equations and returns a three-term certificate; the witness search
reaches candidate 23, $w(x)=x^2+1$, with three candidates reaching universal
certificate search; the Bernstein case $x^2-x+3/10$ on $[0,1]$ needs one split
and yields two leaves each with exact lower bound $1/20$.

The deliberate control is $(x+y+1)^2$: the restricted dictionary finds no
certificate, while a manually supplied one-square certificate is accepted. This
run therefore has a \emph{checker more general than its searcher}, which it uses
explicitly as a guard against an inflated ``SOS-complete'' claim.

This run has the most systematic mutation testing in the design round, with
parametrised
bad-input matrices for cone, box, Horn and induction certificates, including
rejection of a base-case use of the induction hypothesis, invalid generalisation
metadata, rigid-variable intrusion, and reference to an unavailable future
lemma.

\subsection{Four Horn experiments that must not be merged}

Four runs measured demand-directed versus forward Horn search. They used four
incompatible setups, and the temptation to combine their tables should be
resisted.

\begin{table}[htbp]
\centering\small
\begin{tabularx}{\textwidth}{@{}lX@{}}
\toprule
Run & Setup and headline observation \\
\midrule
\src{p5} & Depth sweep 4/8/10/12/14 on $P(a)$, $P(x)\to P(g(x))$,
$P(x)\to P(f(x))$; no distractors; no forward cap. Forward generates
$2^{d+1}-1$ facts (31, 511, 2{,}047, 8{,}191, 32{,}767) against $d+1$ demand
atoms; both proofs have $d+1$ nodes. Timings reported (846.283\,ms against
0.478\,ms at depth 14). \\
\src{p9} & Same rule family, depths 4/8/12/16, but the forward baseline
\emph{stops when the target appears}, giving $2^d$ facts; a 5{,}000-fact cap is
reached at depth 16. No timings reported. \\
\src{p3} & 12 curated goals with decoys placed first; 3/12 unsliced against
12/12 sliced; 250 random graphs and 500 agreement comparisons. \\
\src{p8} & Fixed useful chain of length 12 with a distractor sweep to 1{,}000
chains; 24{,}012 / 11{,}012 / 12 firings for full closure, early-goal stop, and
demand slicing. Timings explicitly disowned. \\
\bottomrule
\end{tabularx}
\caption{The same mechanism, four incomparable measurements. \src{p5} and
\src{p9} differ only in the forward baseline's stopping rule, and that alone
changes the reported growth from $2^{d+1}-1$ to $2^d$.}
\label{tab:horn-four}
\end{table}

That last point is the useful methodological lesson: a factor-of-two difference
in a headline growth curve came entirely from a baseline policy decision, not
from the mechanism under test.

\subsection{The extension round: nine more runs, still not pooled}
\label{sec:extension-runs}

A second round of nine independent proposals, labelled \src{e1} to \src{e9},
was prepared against the merged draft. Their runs are reported here separately
from the first nine and separately from each other. The non-pooling rule applies
with the same force --- and with an additional reason: the two rounds attack
overlapping problems, so a combined total would double-count the same
mathematics under different unit names.

The environmental coincidence repeats almost exactly. All nine ran on Python
3.13.5 on Linux 6.18.44 x86-64 with SymPy 1.14.0, and five of them
(\src{e1}, \src{e2}, \src{e3}, \src{e4}, \src{e6}) recorded the seed
\code{20260915} --- one day after the \code{20260914} shared by five of the
design-round runs. \src{e5} used \code{2026091507}, \src{e9} used
\code{681437}, and \src{e7} and \src{e8} record fixed seeds without a shared
value. As before, the shared seed drives entirely different generators over
entirely different case sets, and the coincidence carries no comparability.

\begin{table}[htbp]
\centering\small
\begin{tabularx}{\textwidth}{@{}llX@{}}
\toprule
Run & Headline unit & Recorded outcome \\
\midrule
\src{e1} & 124 runs / 62 problems &
74 proved, 34 refuted, 16 \unknown{}; 108 accepted certificates, all 108
replayed with site packages disabled and neither the search module nor SymPy
imported (0.052\,s); 40 unit tests. \\
\src{e2} & 65 accepted certificates &
Four families (25 two-sided, 24 invariant, 15 Gosper, 1 telescoper);
1{,}100 assertions; 188 mutations rejected; 527 pointwise binomial regressions
over $n=0,\ldots,30$; 100 differential substitution cases; 7 expected
\unknown{}s. \\
\src{e3} & 317 cases / 313 bundles &
178 proved, 135 refuted, 4 \unknown{}; 240 independent finite-state oracle
comparisons; 210 diagnostic moment evaluations; 25 targeted invalid mutations,
all rejected. \\
\src{e4} & 82 corpus entries / 73 certificates &
Machine lane 37 proved, 28 refuted, 1 \unknown{}; IPC lane 8 countermodels and
8 \unknown{}; 222 generated local Lean obligations, none compiled; largest
certificate 915 bytes; whole corpus 0.842\,s over three repeats; 151 pytest
items. \\
\src{e5} & 59 cases &
38 certified, 17 counterexamples, 4 \unknown{}, across 36 linear, 13 ideal and
10 binomial-sum cases; 818 mutations generated, 0 accepted; 117 pytest items. \\
\src{e6} & 82 attempts &
Ideal lane 27 certificates / 2 \unknown{}; weighted lane 23 closures and 25
distinguishing words; telescoping lane 3 certificates / 2 \unknown{}; 78
certificates replayed under a standard-library-only checker; 318 pytest
items. \\
\src{e7} & 135 replayed records &
113 constant-action invariant, 14 ideal-preservation, 4 ranking, 1
indexing-algebra, 3 orbit-counterexample records; 31 corruptions rejected; 29
test methods; whole run about 1.8\,s. \\
\src{e8} & 422 stored certificates &
51 constant-multiplier, 14 polynomial-multiplier, 136 finite-algebra, 221
integer-projection; all replay with the search entry points replaced by
exceptions; 3{,}400 parameter valuations checked against a complete concrete
oracle (1{,}368 feasible, 2{,}032 infeasible, no disagreement); 21 test methods
containing 19 explicit mutation cases. \\
\src{e9} & 245 accepted objects &
54 matrix closures, 115 separating words, 23 telescopers, 6
common-left-multiple identities, 47 singularity/seed plans; 245 specified
invalidating mutations, all rejected; 1{,}140 exhaustive word evaluations; 483
additional recurrence checks; 41 unit tests; 5.558\,s for the corpus and
0.339\,s for replay with SymPy absent from the replay process. \\
\bottomrule
\end{tabularx}
\caption{The extension round. Nine incommensurable ledgers: ``runs'',
``cases'', ``bundles'', ``certificates'', ``records'', ``attempts'' and
``objects'' are seven different units, and no row may be added to any other.}
\label{tab:extension-runs}
\end{table}

\subsubsection*{The ablations that carry the argument}

Three of these runs contain a controlled comparison that is more informative
than any of the totals, because in each case the \emph{same} problem set is run
through two lanes that differ in exactly one certificate language.

\src{e1} ran sixteen groups: eight problem families, each in a linear-span mode
and an ideal mode. The two modes agree everywhere except in two families. On
\emph{coupled nonlinear} and \emph{nonlinear equivalence}, eight cases each, the
linear lane returns \unknown{} on all sixteen and the ideal lane proves all
sixteen. That is the entire content of Section~\ref{sec:span-insufficient}
reduced to a measurement: no span budget closes those systems, and the failure
is not a fuel problem.

\src{e2} ran 24 invariant cases through the same feature set twice. The subspace
method certifies 24; restricting the certificate language to conserved
quantities certifies 3. The other 21 are true, and the restricted lane simply
cannot say so.

\src{e6} separates its weighted lane's outcomes rather than reporting a success
rate: of 48 attempts, 23 produce closure certificates and 25 produce
distinguishing words. Both are results. Counting only the first would report a
48\,\% success rate for a worker that actually answered every question it was
asked.

\subsubsection*{\src{e7}'s composition table, and how to read it}

One table deserves separate handling because it is the most quotable number in
the extension round and the most easily misquoted. Over four finite hole
grammars held fixed in size and order, a flat whole-term policy and a
contract-directed local policy both return the same term:

\begin{table}[htbp]
\centering\small
\begin{tabular}{@{}rrrrl@{}}
\toprule
Distractors & Flat complete terms & Local checks & Local complete terms & Hole sizes \\
\midrule
0  & 24     & 13 & 1 & 1, 2, 2, 8 \\
4  & 2{,}148  & 29 & 1 & 5, 6, 6, 12 \\
8  & 14{,}384 & 45 & 1 & 9, 10, 10, 16 \\
12 & 50{,}940 & 61 & 1 & 13, 14, 14, 20 \\
\bottomrule
\end{tabular}
\caption{Same grammar, same within-hole ordering, different composition
strategy. Local checks and complete terms are different work units and are not
comparable to each other.}
\label{tab:composition}
\end{table}

The distractor ordering is deliberately unfavourable and is held fixed across
both policies. The largest case was additionally checked on all lists of length
at most six over each of two two-element alphabets, at ordinary nearby indices
and at $\pm2^{80}$, agreeing with an independent non-fold reference on all
3{,}570 checks --- with the symbolic replay, not those observations, carrying
the general claim. This is a demonstration that a cross-product disappears under
factored contracts. It is not a measured $50\,940$-fold speedup, not robustness
under all candidate orders, and not evidence that the original Leant query has
been fixed.

\subsubsection*{What the re-run added}

Unlike the design round, whose authoring environments were unavailable here,
every extension suite could be executed in the merge environment. All nine were
re-run, and every recorded count was reproduced:

\begin{table}[htbp]
\centering\small
\begin{tabular}{@{}llll@{}}
\toprule
Run & Command & Recorded & Observed \\
\midrule
\src{e1} & \code{pytest -q test\_closure.py}    & 40 tests & 40 passed \\
\src{e2} & \code{pytest -q test\_regressions.py}& 13 methods & 13 passed, 69 subtests \\
\src{e3} & \code{pytest -q tests}               & 12 methods & 12 passed, 313 subtests \\
\src{e4} & \code{pytest -q tests}               & 151 tests & 151 passed \\
\src{e5} & \code{pytest -q tests}               & 117 items & 117 passed \\
\src{e6} & \code{pytest -q tests}               & 318 passed & 318 passed \\
\src{e7} & \code{pytest -q test\_extensions.py} & 29 methods & 29 passed, 140 subtests \\
\src{e8} & \code{python -S -m unittest test\_prototypes} & 21 methods & 21 passed \\
\src{e9} & \code{pytest -q tests}               & 41 tests & 41 passed \\
\bottomrule
\end{tabular}
\caption{Re-run in the merge environment --- Python 3.14.4, NumPy 2.4.4, SciPy
1.17.1, SymPy 1.14.0, pytest 9.0.3, Windows. Recorded in
\texttt{results/extension-suites-rerun.json}.}
\label{tab:extension-rerun}
\end{table}

Two observations, neither of them a defect. \src{e8} runs under \code{python -S},
which confirms its standard-library-only claim by making site packages
unreachable. \src{e9}'s suite must be run from \code{prototype/}; from the
package root it fails collection with
\code{ModuleNotFoundError: forge\_summaries}. That is a path convention.

The re-run is a reproduction on \emph{different} interpreter and library
versions than the authoring environment, not a bit-identical replay, and no
timing was compared. It establishes that the suites pass and that the recorded
counts are accurate. It does not establish that any checker is correct, and it
is not a Lean result.

\subsection{The third round: nine more, and a seed that now spans two rounds}
\label{sec:round-three-runs}

A third round of nine proposals was prepared against the merged draft at commit
\code{c98e47c5}, which is to say against the eighteen-proposal document
\emph{including} Section~\ref{sec:closure}. They are labelled \src{r1} to
\src{r9} and reported separately, for the reasons that now apply three times
over.

The seed coincidence has stopped being a coincidence worth explaining and
started being a standing hazard. Four of these nine --- \src{r3}, \src{r4},
\src{r7}, \src{r9} --- recorded the seed \code{20260915}, which is the same
value five of the extension-round runs recorded. The same integer therefore
labels runs in two different rounds, over entirely unrelated generators:
polyhedral and noncommutative and analytic cases in one, exponential-polynomial
ladders in another, Markov decision processes in a third, and pushdown systems
in a fourth. \src{r1} uses seed \emph{ranges} instead
($10000$--$10239$ and $20000$--$20039$), \src{r2} uses \code{2026091503}, and
the remainder record none.

\begin{table}[htbp]
\centering\small
\begin{tabularx}{\textwidth}{@{}llX@{}}
\toprule
Run & Headline unit & Recorded outcome \\
\midrule
\src{r1} & 8 fixture certificates &
240 exhaustive finite-language differential cases agreeing 120 \textsc{proved}
/ 120 \textsc{refuted}; 40 quotient certificates; 20 deliberately invalid
objects rejected; 3 cutoff cases returning \unknown{} with no certificate; 18
test methods. \\
\src{r2} & 119 evidence objects &
90 queries, 20 countermodels, 9 further records; 195 falsifying mutations
rejected over four distinct fault shapes; 57 test methods, one holding 24
atom-permutation subtests; 20 non-receipt Gram queries stored \emph{outside}
the 119. \\
\src{r3} & 239 cases in three lanes &
Polyhedral 110 (79 proved, 30 refuted, 1 \unknown{}); noncommutative 85 (82
proved, 3 \unknown{}); analytic 44 (33 interval, 4 jet, 7 \unknown{}); 342
mutations rejected; 160 vertex-oracle cases and 288 Weyl polynomial actions as
independent checks. \\
\src{r4} & 284 certificates &
25 boundary-zero, 240 random eventual-sign, 8 named, 4 auxiliary-factor, 4
compact-cover, 2 growth-witness, 1 counterexample; an order ablation checking
14{,}409 distinct cofactor orders over 200 inputs with 0 mismatches; 18 test
methods. \\
\src{r5} & 246 records &
62 anchored remainders, 43 ladders, 60 divergence certificates, 81 barriers,
over 230 distinct source/domain subjects; 25 test methods. \\
\src{r6} & 304 records &
196 positive, 40 system, 44 negative point, 8 grammar obstruction, 16 root;
624 targeted invalid mutations rejected; cofactor count reduced from 1{,}212 to
288 on the same 196 problems; 22 test methods. \\
\src{r7} & 520 certificates &
Six workers; 520 mutations rejected, one per certificate; 600 polynomial
finite-horizon checks held separate from the record total; 5 \unknown{}
controls; 24 test methods. \\
\src{r8} & 1{,}564 problems &
884 parity arenas (584 exhaustive at 1--2 vertices, 300 random), 400 transport,
180 MDP, 100 simulation; 2{,}117 mutations; an ablation in which 211 of 400
optimised couplings beat the independent product; 14 regression methods. \\
\src{r9} & 28{,}566 records &
22{,}300 game arenas (22{,}100 exhaustive at 1--3 vertices), 26{,}915 region
records, 67{,}260 start-state replays, 970 Markov chain cases, 881 pushdown
cases; 36 test methods; whole corpus in 1.62\,s. \\
\bottomrule
\end{tabularx}
\caption{The third round. Nine more incommensurable ledgers, and three more
unit names --- ``objects'', ``records'' and ``arenas'' --- none of which is any
of the previous rounds' units.}
\label{tab:round-three-runs}
\end{table}

\subsubsection*{Three traps in this round's own numbers}

Each of these was found by reading the recorded files rather than the prose,
and each would survive a careless summary.

\paragraph{Identical totals over different runs.} \src{r7} ships a pilot
summary and a final summary. Both report \code{records: 520}, both report
\code{rejected\_mutations: 520}, and both record seed \code{20260915}. They are
different runs: the transport lane splits 81/99 in one and 83/97 in the other,
the median bisimulation check time differs, and the maximum certificate size
differs by nearly 300 bytes. A reader comparing headline numbers would conclude
they were the same execution. The proposal says plainly that the pilot ``is not
combined with the final counts'', and this document quotes only the final.

\paragraph{One corpus counted in two files.} \src{r8} ships a smoke run and a
full run. Both report \code{exhaustive\_parity\_arenas: 584}. It is the
\emph{same} 584 arenas --- the exhaustive sweep is identical and only the random
supplement differs, 25 against 300. Adding the two reports gives 1{,}168
exhaustive arenas, and there are 584.

\paragraph{Three units that look like one.} \src{r9} distinguishes a model case,
a start-state check, and a stored record, and the three counts differ because a
single game can yield two nonempty regions and be replayed from many start
states. The numbers are 22{,}300 arenas, 26{,}915 region records and 67{,}260
start-state replays. None of them is the number of theorems, and the proposal
states the distinction as a rule rather than a footnote.

\subsubsection*{Two numbers that are not what they look like}

\src{r9}'s pushdown corpus contains 39 cases in which the bounded reference
oracle \emph{reached its limit and returned no verdict}. They are counted among
the 250 general cases and among the 180 \code{unreachable} records, but the
oracle did not confirm them. The honest decomposition is 70 confirmed
reachable, 141 exhaustively confirmed unreachable, and 39 unsettled; a summary
claiming 250/250 oracle agreement would be a fabrication, and the proposal
records the limitation rather than counting a truncated search as corroboration.

The same proposal's largest number, $2{,}147{,}483{,}647$ steps, is
\emph{represented and not executed}. It was computed from a 31-node DAG;
concrete expansion was checked through depth ten only. \src{r1}'s $2^{60}$ is
the same kind of number arrived at the same way, and both efforts label it as
such.

\subsubsection*{What the re-run added}

Every round-three suite executes in the merge environment and every recorded
count is reproduced --- 18, 57, 12, 18, 25, 22, 24, 14 and 36 test methods for
\src{r1} through \src{r9}. Eight of the nine ship a replay entry point and all
eight return zero, replaying 119, 228, 284, 246, 304, 520, 1{,}564 and 28{,}566
stored objects respectively. \src{r1} has no standalone replay script; its
solver replays each certificate it returns, and its record of eight fixtures
replayed in isolated processes was read rather than re-executed.

Seven of the nine invoke their own suites under \code{python -S}, every replay
ran with site packages disabled, and several block their own search modules
during replay and report having done so. Recorded in
\texttt{results/round-three-suites-rerun.json}. As before, this establishes that
the suites pass, that the recorded counts are accurate, and that the
certificates replay. It establishes nothing about checker correctness.
\subsection{The fourth round: nine more, and one seed in seven of them}
\label{sec:round-four-runs}

A fourth round of nine proposals was prepared against the merged draft at
commit \code{58ea206}, which is to say against the twenty-seven-proposal
document including Sections~\ref{sec:noncommutative} to~\ref{sec:quantitative}.
They are labelled \src{s1} to \src{s9}.

The seed observation has now completed its arc. In the design round five runs
shared \code{20260914}; in the extension round five shared \code{20260915}; in
the third round four shared that same value; and in this round \emph{seven of
the nine} record \code{20260915} --- \src{s1}, \src{s3}, \src{s5}, \src{s6},
\src{s7}, \src{s8}, \src{s9}. Across three rounds, sixteen runs now carry one
integer, over sixteen unrelated generators: polyhedral cells, exponential
ladders, Markov decision processes, pushdown systems, Petri nets, real-root
covers, binary automata, permutation groups and chain complexes. One of them
anticipates the hazard and defuses it in its own text --- ``these seeds have
meaning only together with the supplied generators and are unrelated to
similarly numbered experiments in Forge'' --- and another says the same. The
rest do not. A matching seed in this collection is not evidence of a shared
corpus, and by now it is barely evidence of anything.

All nine also record byte-identical Python and platform strings, which means
they ran on one host image. That has a consequence worth stating: the
\code{NOT\_RUN} Lean status appearing nine times is \emph{one} environmental
fact reported nine times, not nine independent findings.

\begin{table}[htbp]
\centering\small
\begin{tabularx}{\textwidth}{@{}llX@{}}
\toprule
Run & Headline unit & Recorded outcome \\
\midrule
\src{s1} & 423 replay records &
94 coverability frontiers, 136 least unsafe parameters, 69 parameter
frontiers, 63 equal and 58 different guarded endpoint equivalences, plus
compressed runs and lassos; 68 unit tests; 8 invalid records rejected; 4
cutoff controls returning \unknown{}. \\
\src{s2} & 2{,}360 receipts &
2{,}100 exhaustive Petri instances (1{,}536 safe, 564 coverable), 120 random,
45 scalar and 54 coupled probabilistic systems, 26 spectral extinction cases,
8 algebraic decoders; 275 rejection controls; no unit-test suite. \\
\src{s3} & 650 stored objects &
305 counter-system, 163 nominal, 182 mixed; 27 test methods; one Lean specimen
that elaborates. \\
\src{s4} & 190 and 193 certificates &
Eager and separation producers over the same models, agreeing on every basis;
1{,}296 local predecessor problems and 63{,}504 membership comparisons against
an oracle; 11{,}520 complete finite execution queries; 37 corruption attempts
rejected; 4 cutoff controls. \\
\src{s5} & 100 replayed atlases &
24 primary examples with 356 plane cells, 76 metamorphic derivations, 250 root
problems, 200 localised trials; 540 mutations; 42 pytest items. \\
\src{s6} & 99 valid bundles &
748 exact count comparisons over 19 quantified problems and 113 cells; 973
invalid mutations rejected; 160 Boolean formulas and 4{,}800 truth
assignments; 25 unit tests. \\
\src{s7} & 90 certificates &
25 named queries, 60 guarded alternating formulas, 5 witness-relation graphs;
643 concrete witnesses including one at $10^{100}+37$; 1{,}096 mutation
objects; 7 interface tests. \\
\src{s8} & 8{,}028 membership queries &
3{,}852 positive and 4{,}176 negative; 160 chain bundles; 296 canonical images
over 24{,}710 cover nodes; 40 cycle inventories; 33 corruption controls; 12
test methods; $40!$ certified in 11{,}975 bytes. \\
\src{s9} & 656 certificates &
97 homology coordinates, 129 fillers, 65 obstructions, 48 induced maps, 72
reductions, 100 homotopies, 145 no-homotopy obstructions; 489 rejection
controls; 12 test methods. \\
\bottomrule
\end{tabularx}
\caption{The fourth round. Nine more incommensurable ledgers, and the units
have stopped being nameable in common: ``records'', ``receipts'', ``objects'',
``certificates'', ``atlases'', ``bundles'' and ``queries'' are seven different
things here, and none of them is a theorem.}
\label{tab:round-four-runs}
\end{table}

\subsubsection*{The traps this round contains}

Three are of a kind already seen, and two are new.

\paragraph{Superseded runs that reproduce their own totals.} Two proposals ship
two runs each and both say, in their own words, not to add them. One records
``96 valid bundles and 944 rejected mutations'' in an earlier run and warns:
``Do not add these numbers to run-02. Most mathematical cases are shared.'' The
other: ``Runs 1 and 2 use the same seed, generators and mathematical cases.
Their counts must not be summed or represented as independent coverage.'' This
is the same shape as the previous round's identical-headline pilot, and it is
now the third round in which it has occurred.

\paragraph{A number that appears twice inside one proposal.} \src{s7}'s 643 is
both its count of accepted witness traces and the size of one mutation class
inside its 1{,}096 total. The mutation table sums to exactly 1{,}096 with that
class included, so quoting ``643 witnesses and 1{,}096 mutations'' is correct
while adding them is not.

\paragraph{Twelve and twelve.} \src{s8} and \src{s9} each report exactly twelve
unit-test methods. The two lists share no name and no subject --- one tests
alternating parity and noncommuting composition, the other tests transposition
and rectangular multiplication. Numerically identical, entirely unrelated, and
the single most likely accidental conflation in the collection.

\paragraph{Three impossibility results are three.} \src{s5}, \src{s6} and
\src{s7} each prove a small impossibility, by three different techniques, with
one shared moral. They are recorded separately in
Section~\ref{sec:realalgebraic}; a ledger citing them as one confirmation of
one fact would be wrong three times over.

\paragraph{A common-mode risk no single proposal could see.} \src{s5} and
\src{s6} both use the same computer-algebra package as their untrusted
producer, and both ran in this round's shared environment. Each discloses the
dependency \emph{within itself}. Neither is in a position to observe that a
defect in that package would be correlated across two independent proposals,
and this document is.

\subsubsection*{What the re-run added}

Every round-four suite that exists runs in the merge environment and passes:
68, 27, 42, 25, 7, 12 and 12 test methods for \src{s1}, \src{s3}, \src{s5},
\src{s6}, \src{s7}, \src{s8} and \src{s9}, plus \src{s4}'s combined
\code{run\_tests.py} reporting \code{REPLAY\_OK}. \src{s2} ships no unit suite
at all; its evidence is the replay corpus and in-experiment assertions.

Seven replay entry points return zero, covering 423, 2{,}360, 650, 99, 90,
8{,}028-query and 656 stored objects respectively. Recorded in
\texttt{results/round-four-suites-rerun.json}.

\begin{remark}[Two proposals, opposite instructions about \code{-O}]
\src{s2} instructs: \emph{do not} run under \code{python -O}, because its
experiment assertions must remain enabled. The previous round's \src{r7} ran
its suite under \code{-O} \emph{specifically} so that no check could be
implemented as a removable \code{assert}. Both were followed as written here.
They are opposite conclusions from the same observation --- that \code{assert}
disappears under optimisation --- and the second is the safer discipline: a
check that vanishes under a flag was never a check.
\end{remark}

\subsubsection*{Three lanes merged into the prototype}

Three of this round's four antichain lanes are folded into the merged package
as \code{forge/wsts/}: \src{s1}'s frontiers with compressed witness runs,
\src{s3}'s finite-control counter systems with parameterised initial families,
and \src{s4}'s matrix updates and lossy channels. All three are standard
library only with their searches included, so they preserve the property that
makes the package's replay evidence worth anything: the whole of it runs under
\code{python -S}.

The merge is where the duplication became visible as code rather than as prose.
The four antichain proposals wrote one well-quasi-order, one antichain, one
predecessor formula and one saturation loop three or four times between them.
The merged package states each once, and a single backward loop takes the
predecessor rule as an argument, so it does not know which of the four model
languages it is driving. What still appears four times is what genuinely
differs: the model languages themselves.

Two distinctions were preserved rather than flattened. \src{s1}'s producer
composes run summaries in $(	ext{need},	ext{give})$ coordinates and its
checker recomputes them in $(	ext{need},	ext{signed delta})$ coordinates;
the two recurrences look different and must agree through
$	ext{give}=	ext{need}+	ext{delta}$, which is the strongest
producer/checker separation this round contains, and it survives the merge with
a test asserting exactly that identity. And a matrix-update transition has no
unique least predecessor, so the merged interface returns a list where a vector
addition would return one element --- collapsing that would have made the
backward search unsound rather than merely incomplete.

The package's own numbers, which are its own and are not a fraction of anything:
776 tests pass, 35 stored certificates recheck, and 53 mutations are rejected.
Before this merge the suite reported 748, and \src{s6} reports 748
specialization evaluations; the two were never related, and the collision
closing by accident is a reminder that a repeated number in this collection is
not evidence of a shared quantity.

Five lanes are deliberately not merged: \src{s2}'s probabilistic fixed points,
which share the word ``coverability'' with \src{s1} and nothing else;
\src{s3}'s equality-register transducers, which are finite by the Bell numbers
rather than by any well-quasi-order; and \src{s5} and \src{s6}'s real cell
covers, which need a computer-algebra producer. The third round remains
unmerged in full, for the reasons given in \S
ef{sec:noncommutative}, and in
its analytic and quantitative sections.

\subsection{Lean status}

\subsubsection*{What the thirty-six runs recorded}

None of the thirty-six compiled any Lean source. Every one of them recorded the
same reason --- no \code{lean} or \code{lake} executable in the authoring
environment --- under several different filenames and in more than one format.
Consequently none performed a kernel check, an axiom audit, or any comparison
against \grind{}, Aesop, \code{nlinarith}, LeanHammer, or any other tactic.

The quantity of candidate Lean source, by round, is the most legible trend in
this document:

\begin{center}\small
\begin{tabular}{@{}lrr@{}}
\toprule
Round & Source files & Lines \\
\midrule
Design \src{p1}--\src{p9}     & 21 & 3{,}022 \\
Extension \src{e1}--\src{e9}  & 20 & 3{,}144 \\
Third \src{r1}--\src{r9}      & 16 & \phantom{0}755 \\
Fourth \src{s1}--\src{s9}     & \phantom{0}2 & \phantom{00}98 \\
\bottomrule
\end{tabular}
\end{center}

It did not decline because the proposals got smaller. It declined because one
of them argued that an uncompiled file is not evidence, and a compiler then
agreed.

\src{r9} was first: it ships no Lean at all, and says why --- it ``deliberately
contains no placeholder theorem files presented as implemented proofs''. That
is a pointed remark about the practice the other twenty-six follow, and the
next subsection shows it landing. By the fourth round, four proposals ship no
\code{lean/} directory at all and three ship one containing only an obligation
document, three of them in near-identical words: the directory ``deliberately
contains no \code{sorry}-based or uncompiled theorem file presented as a
completed implementation''.

\subsubsection*{What this merge was able to add}

The merge environment has \code{elan}, and installing the pinned
\code{leanprover/lean4:v4.34.0} made a first, narrow compilation check possible.
Three of the design-round files in the merged tree import nothing beyond Lean
core and can therefore be elaborated without a built Mathlib:

\begin{table}[htbp]
\centering\small
\begin{tabular}{@{}llr@{}}
\toprule
File & Result & Elapsed (s) \\
\midrule
\src{p3} \code{lean/DesignAPI.lean}   & elaborated, no errors & 211.4 \\
\src{p6} \code{lean/Contracts.lean}   & elaborated, no errors & 5.3 \\
\src{p9} \code{lean/Structural.lean}  & elaborated, no errors & 5.1 \\
\bottomrule
\end{tabular}
\caption{Lean v4.34.0 (commit \texttt{293d5d0c}), recorded in
\texttt{results/lean-core-elaboration.json}.}
\label{tab:lean}
\end{table}

This confirms that the proposed runtime and scheduler contract, the
\code{CertificateSpec} soundness contract, and the Mathlib-free structural list
proofs are well-typed Lean 4.34.0 rather than merely plausible-looking source.

It confirms nothing else. It is not a transitive axiom audit. It says nothing
about the thirteen files in the merged tree that import Mathlib and need a
built project, or about the two that merely import their siblings. It does not establish that any certificate checker is correct.
And it is not a comparison against any Lean tactic --- successful elaboration of
a hand-written example proves that the example type-checks, not that automation
found it.

\subsubsection*{Ten more files, from the extension round}

The same check was then applied to the extension proposals. Of their twenty
source files, two are aggregators that import their siblings, eight import
Mathlib, and ten import nothing beyond Lean core. All ten elaborate:

\begin{table}[htbp]
\centering\small
\begin{tabular}{@{}llr@{}}
\toprule
File & Result & Elapsed (s) \\
\midrule
\src{e1} \code{lean/Core.lean}                          & elaborated & 82.1 \\
\src{e2} \code{lean/ForgeExtensions/Reachability.lean}  & elaborated & 175.8 \\
\src{e3} \code{lean/BridgeSpec.lean}                    & elaborated & 220.7 \\
\src{e4} \code{lean/Reachability.lean}                  & elaborated & 14.7 \\
\src{e5} \code{lean/ForgeExtension/Core.lean}           & elaborated & 9.4 \\
\src{e6} \code{lean/ProofPrinciples.lean}               & elaborated & 10.2 \\
\src{e7} \code{lean/Specimens.lean}                     & elaborated & 9.5 \\
\src{e8} \code{lean/CoreSoundness.lean}                 & elaborated & 8.8 \\
\src{e8} \code{lean/IndexingSketch.lean}                & elaborated & 30.8 \\
\src{e9} \code{lean/WordSimulation.lean}                & elaborated & 2.2 \\
\bottomrule
\end{tabular}
\caption{Lean v4.34.0 (commit \texttt{293d5d0c}), recorded in
\texttt{results/lean-extensions-elaboration.json}.}
\label{tab:lean-extensions}
\end{table}

Seven of the ten additionally print \code{\#print axioms} results, and every
printed line reads \emph{does not depend on any axioms} --- covering the
reachability and run invariants of \src{e2}, \src{e4}, \src{e6} and
\src{e8}, the tree-cover lemmas of \src{e8}, the fold-indexing and ranked-call
lemmas of \src{e7}, and the word-simulation lemmas of \src{e9}. This is a
stronger statement than the design-round result, where two of the three
elaborated theorems depend on \code{propext}.

\subsubsection*{The third round, and the first file that did not compile}

Of the third round's sixteen Lean sources, thirteen import Mathlib and one
proposal ships none. That leaves three, and for the first time in the
collection's history a compiler had something to say about delivered Lean.

\begin{table}[htbp]
\centering\small
\begin{tabular}{@{}llr@{}}
\toprule
File & Result & Elapsed (s) \\
\midrule
\src{r3} \code{lean/CoreComposition.lean}     & elaborated & 34.1 \\
\src{r7} \code{lean/ForgeQ/Monotone.lean}     & elaborated; no axiom dependencies & 106.7 \\
\src{r8} \code{lean/RankTelescoping.lean}     & \textbf{failed}, ten errors & 3.6 \\
\bottomrule
\end{tabular}
\caption{Lean v4.34.0, recorded in
\texttt{results/lean-round-three-elaboration.json}.}
\label{tab:lean-round-three}
\end{table}

The failure is instructive precisely because it is trivial. \src{r8}'s file
defines \code{prefix}, which is a reserved command keyword in Lean~4. All ten
errors cascade from that single identifier, ending in \emph{unknown constant}
for the theorem below it, which never parsed. The mathematics is fine: renaming
the definition on a scratch copy makes both theorems elaborate, depending on
\code{propext} and \code{Quot.sound} by way of \code{omega} and \code{simp}.

\begin{remark}[The broken file prints a clean axiom line]
The last line \src{r8}'s file emits is
\code{'ForgeAP.prefix\_bounded' does not depend on any axioms}.

It exits~1 with ten errors above that line. Lean recovers from parse errors and
keeps processing, so the \code{\#print axioms} command attached to a
declaration that \emph{did} elaborate still ran and still printed. A reader who
looks at the tail of that output, or any tool that greps for the phrase without
checking the exit code, reads this file as clean.

This document's own sweep script made exactly that mistake and counted thirteen
files reporting no axiom dependencies where twelve elaborated --- which is how
the case was found. It is the second tooling defect this merge discovered in
itself, after the \code{sorry} scan, and it has the same shape: a substring
taken as a verdict. The figure quoted everywhere in this document is twelve,
counting only files that also compiled.

The lesson is narrower than ``do not trust \code{\#print axioms}''. The
command is sound; it reports on the declaration it was given. What carries no
evidence is the \emph{line}, detached from the exit status of the file that
produced it --- which is the form in which such lines usually travel.
\end{remark}

Nothing about this is a mathematical criticism of that proposal, and the repair
was deliberately \emph{not} applied to the archived source --- the proposals
are frozen as delivered, and the fix belongs in the merged tree if the lemma is
adopted. What it demonstrates is the thing \src{r9} asserted and twenty-six
proposals had no way to discover: an uncompiled \code{.lean} file with an
unexecuted \code{\#print axioms} command at the bottom carries exactly as much
evidence as its author's care, and no more. Two of these three carried that
care. One did not, and only a compiler could tell them apart.

\subsubsection*{A second broken file, found only by sweeping everything}

The table above covers the three core-only files, which is the set every
round's scan looked at. A final sweep over \emph{all} eighty-six
\code{.lean} files in the repository --- Mathlib-importing ones included, on
the expectation that they would fail for want of Mathlib and could be dismissed
as a class --- found one that fails for a different reason.

\src{r6}'s \code{lean/FlowTargets.lean} opens with a module docstring and
places \code{import Mathlib} on line~9, below it. In Lean~4 an \code{import}
must precede everything, including a docstring, so the file stops at
\code{invalid 'import' command, it must be used in the beginning of the file}.
The error is at parse time and before import resolution, so Mathlib's absence
is not the cause: a four-line file consisting of a docstring, \code{import
Init} and a trivial theorem reproduces it exactly, and moving the import to the
top makes that file elaborate. \src{r6}'s would parse under the same repair.

So the count of delivered Lean a compiler has contradicted is \emph{two}, not
one, and the two are different failures. \src{r8}'s is a core-only file that
was expected to elaborate and did not. \src{r6}'s is a file nobody had ever
tried, whose first four lines are wrong in a way that has nothing to do with
its mathematics --- and whose own header states that it ``was not compiled in
the creation environment,'' which is true, and which is why nobody found out.

The asymmetry is the point. Forty-eight files import Mathlib and eleven import
sibling modules; none of the fifty-nine has ever been checked by anything, and
the one time this document looked inside that bucket for a reason unrelated to
Mathlib, it found a file that cannot parse. That is one observation and not a
basis for estimating how many others there are. It does establish that
``fails because Mathlib is missing'' was an assumption about those files rather
than a measurement of them.

\begin{remark}[The two rounds also duplicated a lemma in Lean]
\src{r7}'s \code{coupling\_left\_expectation} and \src{r8}'s
\code{left\_expectation} are the same lemma with the same proof script --- a
\code{Finset.sum\_congr} followed by a rewrite through the marginal --- stated
independently. It is the only literal duplication of Lean \emph{content}
across the third round, and it would be merged the way
Section~\ref{sec:closure}'s principles were.
\end{remark}

\subsubsection*{The fourth round, and a defect in this repository's own tools}

Round four ships two Lean files and both elaborate: \src{s3}'s counter lemmas
in 93\,s, depending on \code{propext}, \code{Classical.choice} and
\code{Quot.sound}; and \src{s5}'s atlas-quantifier composition lemmas in
5\,s, with no axiom dependencies for either theorem it exposes.

The first run reported the former as \code{elaborated\_with\_sorry}, and that
was wrong. The scan in \code{tools/check\_lean.py} tested for the substring
\code{sorry} in the raw file. That file's only occurrence of the token is its
own header sentence, \emph{``There are no sorry/admit placeholders''} --- so
the scan reported a defect on a file that was declaring its own cleanliness.
The harness now strips Lean comments before the search, with a depth counter
because Lean block comments nest, and a token found only in prose is recorded
as such rather than as a defect.

Re-auditing all eighty-six Lean files across the four rounds and the merged
tree with the corrected scanner: \emph{none} contains a real \code{sorry} or
\code{sorryAx}. All twenty-five occurrences of the token are authors stating
that there are none. That is a better result than anyone claimed, and it could
not have been established before, because the unfixed scanner could not tell
the two cases apart.

The defect is recorded here rather than quietly fixed because the earlier
rounds' scans were produced by the unfixed version. They happened to be correct
--- none of the files they examined contained the token at all --- but that was
luck, and a tool that false-positives on a comment would equally have missed a
\code{sorry} placed where comment-stripping changes the answer.

The same three disclaimers apply, and one more besides. \src{e8}'s
\code{IndexingSketch.lean} elaborates, and its author had marked it
\emph{uncompiled} and \emph{not a substitute for the original query}. Both
statements are now true at once: the file type-checks, and it still proves an
equation between two ordinary-list definitions rather than anything about
Church-encoded inputs. Elaborating a specimen does not upgrade what the specimen
says.

\subsubsection*{Three more, written by the merge}

Those ten files overlap heavily. Between them they contain six spellings of one
reachability-invariant theorem and five of one word-fold preservation theorem,
differing in constructor names, argument order, and whether the initial state is
a predicate or a point. Merging them means proving each once and keeping what is
actually distinct, so the merged tree gains three files:

\begin{table}[htbp]
\centering\small
\begin{tabularx}{\textwidth}{@{}lXr@{}}
\toprule
File & Merged from & Elapsed (s) \\
\midrule
\code{Forge/Closure/Principles.lean} &
all nine, for the reachability and word-fold theorems; \src{e3} and \src{e4}
for the two-machine relation; \src{e3}, \src{e6} and \src{e9} for the
commuting step map; \src{e1} for the left fold; \src{e7} for well-founded
descent & 6.6 \\
\code{Forge/Closure/Covers.lean} & \src{e8}, the tree half of its soundness
specimen & 244.4 \\
\code{Forge/Closure/Indexing.lean} & \src{e7}'s Church half and \src{e8}'s
\code{List}/\code{Int} sketch, which proved the same equation under different
names & 3.8 \\
\bottomrule
\end{tabularx}
\caption{Recorded in
\texttt{results/lean-closure-elaboration.json}. All three elaborate, and all
twelve of their theorems print \emph{does not depend on any axioms}.}
\label{tab:lean-closure}
\end{table}

Twenty-three files therefore elaborate: six in the merged tree, three in the
design-round proposals, ten in the extension proposals, two of the third
round's three and both of the fourth round's two. Twelve of the twenty-three
report no axiom dependencies. The merged \code{Indexing.lean} inherits
\src{e8}'s scope note verbatim rather than quietly widening it.

The three design-round files are a late correction, and the reason they were
missed is worth stating. Each round's scan looked at the merged tree and at
that round's own proposals; no round went back over \src{p1}--\src{p9}'s own
\code{lean/} directories. \src{p3}'s \code{DesignAPI.lean}, \src{p6}'s
\code{Contracts.lean} and \src{p9}'s \code{Structural.lean} import only
\code{Lean}, and all three elaborate. All three also carry a header declaring
that they were not compiled --- \src{p9}'s says so in capitals --- which is
true of their authoring environment and false of this one.

That is twenty-three of the twenty-four files in this repository that import
nothing beyond Lean core. The twenty-fourth is \src{r8}'s, below. Forty-eight
more import Mathlib directly and eleven import sibling modules; nothing has
ever checked any of those fifty-nine.

\subsection{What the evidence establishes}

It establishes that the implementations actually synthesise certificates in the
stated fragments; that certificate replay is separable from the numerical and
symbolic search libraries; that the specified mutation and corruption tests are
rejected; that several mechanisms solve problems their own restricted precursors
cannot; that all nine extension suites re-run and reproduce their recorded
counts on a different interpreter and different library versions; that all
nine third-round suites and every round-four suite do the same and that their
stored objects replay; and that twenty-three Lean artefacts --- three design
files in the merged tree, three more in the design-round proposals, ten from
the extension round, three the merge wrote by deduplicating those ten, two from
the third round and two from the fourth --- type-check under the pinned
toolchain, twelve of them printing \emph{does not depend on any axioms} for
every theorem they expose.

It also establishes, for the first time, that a delivered Lean file did
\emph{not} type-check --- and, separately, that a scan in this repository's own
tooling reported a defect that was not there.

It does not establish that \forge{} is implemented as a Lean tactic, that the
Mathlib-dependent candidate scripts compile, that any certificate is checked by
Lean's kernel, that the methods outperform existing Lean tactics, or that any
success rate transfers to user developments. Those claims require different
experiments, described next.

Two structural limitations of the evidence deserve to be named rather than
buried. Manufactured examples favour the grammar that generated them --- most
visibly for \src{p2}'s 40 square certificates and \src{p5}'s 100 generated cone
cases, both of which are drawn from the very dictionary their search covers. And
search and replay share substrate: the same sparse polynomial and term
representations appear on both sides in every prototype. Search-free replay is a
meaningful independence property, but it is not an independently implemented
formal verification of that representation. The strongest partial answers to
this are the checkers that run algorithmically \emph{opposite} to their
producers --- expanding a claimed Bernstein basis rather than recomputing
coefficients, re-deriving a clause by independent unit propagation rather than
replaying a resolution chain, recomputing a Sturm chain rather than trusting a
stored root count, verifying $\det U=\pm1$ rather than calling the elimination
that produced $U$.

The extension round sharpens that last point into a three-level scale rather
than a yes/no property, and it is worth stating explicitly because the levels
are easy to conflate. At the strongest level the checker performs a genuinely
different computation from the search: \src{e7}'s ideal lane does not verify a
Gr\"obner basis, re-run Buchberger, or trust a zero-remainder flag, and
\src{e9}'s checkers multiply rational matrices and compare, where the searches
solved nullspaces. At an intermediate level the search entry points are
\emph{replaced by exceptions} during replay --- \src{e8}'s invariant and finite
lanes --- which proves separation but leaves input-table semantics shared. At
the weakest level, \src{e8}'s integer projection reconstructs the witness
program with the same assembler the search used, so replay catches a modified
receipt but not an arithmetic mistake common to both. The proposal says so
itself and mitigates it by a different route, an exhaustive concrete oracle over
a provably sufficient range. A reader comparing ``all certificates replay''
claims across the thirty-six prototypes should ask which of these three things
the sentence means.

One more limitation is specific to the second round and does not appear in the
first. Several extension corpora are \emph{constructed positives}: \src{e7}'s
twelve power-curve fixtures have known invariants, \src{e9}'s forty comparison
models are invertible coordinate changes of random systems, and \src{e8}'s
forty-eight affine-diagonal instances are built to preserve the diagonal. These
exercise polynomial multiplication, substitution and coefficient handling. They
are not held-out benchmarks of discovery performance, and the proposals label
them as such rather than counting them as solved theorems. Likewise \src{e7}'s
seventy-two affine differential systems include forty-eight with a zero relation
space; those are useful controls, not additional nontrivial results.
